\documentclass[11pt]{article}
\usepackage[final]{acl}
\usepackage{times}
\usepackage{latexsym}
\usepackage{fontspec}
\usepackage{polyglossia}
\setdefaultlanguage{russian}
\setotherlanguage{english}
\setmainfont{Times New Roman}
\newfontfamily\cyrillicfont{Times New Roman}
\usepackage{microtype}
\usepackage{inconsolata}
\usepackage{graphicx}

\title{Обучение ранжированию для задач информационного поиска}

\author{Игорь Пивнев, ИТМО}

\begin{document}
\maketitle

\section{Постановка задачи и общий подход}

В данной работе исследуется подход Learning to Rank для задачи информационного
поиска. В отличие от классических моделей ранжирования (например, BM25), которые
используют фиксированную формулу, методы обучения ранжированию позволяют обучать
функцию релевантности на размеченных данных.

В качестве основной модели используется градиентный бустинг над решающими
деревьями, реализованный в библиотеке CatBoost. Обучение проводится в постановке
группового ранжирования (group-wise), где объекты объединяются по запросам, а
целевая функция оптимизирует порядок документов внутри каждой группы.

Качество ранжирования оценивается с помощью метрики nDCG, которая учитывает как
релевантность документов, так и их позиции в ранжированном списке.

\section{Эксперименты на MS LETOR}

Для проверки корректности подхода была использована стандартная коллекция
MSRank-10K из набора LETOR. Данный датасет содержит заранее подготовленные
признаки для пар запрос-документ и оценки релевантности.

Были обучены модели CatBoost с различными функциями потерь:
\begin{itemize}
    \item YetiRank (listwise подход);
    \item PairLogit (pairwise подход).
\end{itemize}

В качестве baseline использовалось случайное ранжирование.

Полученные результаты:
\begin{itemize}
    \item случайное ранжирование: \textbf{0.6336};
    \item YetiRank: \textbf{0.7480};
    \item PairLogit: \textbf{0.7334}.
\end{itemize}

Обе обученные модели существенно превосходят случайный baseline, что подтверждает
корректность применения CatBoost к задаче ранжирования.

\section{Internet Mathematics 2009}

Датасет Internet Mathematics 2009 был предварительно приведён к формату LETOR.
Для каждой строки были извлечены признаки, а документы сгруппированы по запросам.

Основные характеристики набора:
\begin{itemize}
    \item количество объектов: \textbf{77714};
    \item количество признаков: \textbf{245};
    \item количество запросов: \textbf{24425};
    \item среднее число документов на запрос: \textbf{3.1817}.
\end{itemize}

Аналогично предыдущему эксперименту были обучены модели CatBoost:

\begin{itemize}
    \item Random baseline: \textbf{0.8017};
    \item YetiRank: \textbf{0.8985};
    \item PairLogit: \textbf{0.8886}.
\end{itemize}

Результаты показывают, что методы обучения ранжированию устойчиво работают и на
новых данных, даже при отсутствии ручной настройки признаков.

\section{Улучшение BM25 на WikIR}

В данной части работы рассматривалась задача улучшения ранжирования BM25 с помощью
Learning to Rank.

Сначала был реализован baseline на основе BM25, а далее для каждой пары
запрос-документ были извлечены дополнительные признаки:
\begin{itemize}
    \item суммарная частота терминов запроса (TF);
    \item нормализованная частота;
    \item длина документа и запроса;
    \item число общих терминов;
    \item суммарный IDF;
    \item расстояние между термами запроса в документе (proximity).
\end{itemize}

Для обучения использовались все релевантные документы и равное количество случайно
выбранных нерелевантных документов для каждого запроса.

После обучения модели CatBoost она применялась для переранжирования top-100
документов, полученных BM25.

Результаты:
\begin{itemize}
    \item BM25: \textbf{0.8492};
    \item BM25 + LTR (CatBoost): \textbf{0.8473}.
\end{itemize}

Наблюдается ухудшение качества, что подтверждает, что добавление признаков и
обучение модели не всегда позволяет учитывать более сложные зависимости, чем BM25.

\subsection{Реконструкция BM25}

Дополнительно была предпринята попытка аппроксимировать BM25 с помощью обучения
модели на базовых признаках:
\begin{itemize}
    \item TF;
    \item IDF;
    \item длина документа.
\end{itemize}

Несмотря на простоту признаков, обученная модель способна частично воспроизвести
поведение BM25, что подтверждает, что данная формула может быть интерпретирована
как частный случай обучаемой модели.

\section{Эксперименты на MIRAGE}

Для датасета MIRAGE была проведена собственная подготовка данных:
\begin{itemize}
    \item запросы разбиты на train/test по группам;
    \item каждый документ представлен как два поля: заголовок и основной текст;
    \item сформированы признаки на основе совпадений терминов.
\end{itemize}

Используемые признаки включают:
\begin{itemize}
    \item TF в тексте и заголовке;
    \item длины текста и заголовка;
    \item пересечение терминов запроса с текстом и заголовком.
\end{itemize}

После обучения модели CatBoost были получены следующие результаты:
\begin{itemize}
    \item nDCG: \textbf{0.7614}.
\end{itemize}

\section{Выводы}

Проведённые эксперименты позволяют сделать следующие выводы:

\begin{itemize}
    \item методы Learning to Rank существенно превосходят случайные и базовые подходы;
    \item CatBoost эффективно решает задачу ранжирования благодаря поддержке групповых данных;
    \item качество модели сильно зависит от набора признаков;
    \item даже простые признаки (TF, IDF, длина) дают заметный эффект.
\end{itemize}

Таким образом, обучение ранжированию является мощным инструментом для улучшения
качества информационного поиска и может быть использовано поверх классических
моделей, таких как BM25.

Этот и другие отчеты, а также исходный код доступны в репозитории: \url{https://github.com/qw1zzard/itmo-ir-search}.

\end{document}
